\chapter{REFINING THE NOTION OF PROBABILITY}

\textit{July 21st.}

Consider the random walk on the line; a drunk ant sits at the origin and takes steps of length $1$ to the left or right with equal probability. We wish to obtain the probability that after an even $n$ steps, the ant returns to the origin. Of course, the natural approach is to frame this problem in a discrete probability space.

\begin{definition}
    A \eax{discrete probability space} is a pair $(\Omega,q)$ where $\Omega$ is a non-empty finite or countably infinite set called the sample space, and $q:\Omega\to[0,1]$ is a function satisfying $\sum_{\omega\in\Omega}q(\omega)=1$. Any subset of $\Omega$ is termed an event, and the probability of an event $A \subseteq \Omega$ is defined as $\Prob(A) = \sum_{\omega\in A}q(\omega)$.
\end{definition}

For the problem above, we can take the sample space to be
\begin{align}
    \Omega = \{\bv{x}=(x_{0}=0,x_{1},x_{2},\ldots,x_{n}) \in \Z^{n+1} \mid x_{i+1}-x_{i} \in \{-1,1\} \text{ for all } 0 \leq i < n\},
\end{align}
and the mapping $q$ to be $q(\bv{x}) = 2^{-n}$ for all $\bv{x} \in \Omega$. Then, the event that the ant returns to the origin after $n$ steps is given by $A = \{\bv{x} \in \Omega \mid x_{n} = 0\}$, and we have
\begin{align}
    \Prob(A) = \sum_{\bv{x} \in A} q(\bv{x}) = 2^{-n} \abs{A} = 2^{-n} \binom{n}{n/2}.
\end{align}
Reframing the problem, let us consider a random walk on the line starting at the origin, and continuing indefinitely. In this case, the sample space is given by
\begin{align}
    \Omega = \{\bv{x}=(x_{0}=0,x_{1},x_{2},\ldots) \in \Z^{\N} \mid x_{i+1}-x_{i} \in \{-1,1\} \text{ for all } i \in \N\}.
\end{align}
This $\Omega$ is uncountably infinite. In such a case, we cannot define a probability measure by assigning a positive probability to each element of $\Omega$. As another example of difficulty in dealing with uncountable sample spaces, consider the problem of choosing a real number uniformly at random from the interval $[0,1]$. If we attempt to assign a positive probability to each element of $[0,1]$, we would have to assign a probability of $0$ to each element, since there are uncountably many elements in $[0,1]$. However, this would lead to a total probability of $0$, which is not valid. Of course, we know that the probability of choosing a specific element or a countable set of elements from $[0,1]$ is $0$, but there then arises the question of how to assign a probability to an uncountable set of elements, such as the Cantor set. We can work in a slightly different viewpoint.

\begin{definition}
    A probability space is a tuple $(\Omega,\cF,\Prob)$ where $\Omega$ is a non-empty set called the sample space, $\cF$ is the collection of all subsets of $\Omega$ ($\sim 2^{\Omega}$), and $\Prob:\cF\to[0,1]$ is a function satisfying $\Prob(\Omega)=1$ and countable additivity, that is, for any countable collection of pairwise disjoint sets $A_{1},A_{2},\ldots \in \cF$, we have
    \begin{align}
        \Prob\left(\bigcup_{i=1}^{\infty}A_{i}\right) = \sum_{i=1}^{\infty}\Prob(A_{i}).
    \end{align}
\end{definition}

Using this definition, we look to define a $\Prob$ on subsets of $[0,1]$ that satisfies the above properties, along with the property that $\Prob([a,b]) = b-a$. Thus, we have $\Prob([0,1]) = 1$. If we have two disjoint intervals $I_{1},I_{2} \subseteq [0,1]$, then $\Prob(I_{1}\cup I_{2}) = \Prob(I_{1}) + \Prob(I_{2})$. This is a natural property to expect from a probability measure. What about a general set $S \subseteq [0,1]$? One can observe that if we have a sequence of intervals $I_{1},I_{2},\ldots$ such that $S \subseteq \bigcup_{k=1}^{\infty}I_{k}$, then we would have $\Prob(S) \leq \sum_{k=1}^{\infty}\Prob(I_{k})$. This is because the probability of $S$ should be less than or equal to the probability of any set that contains $S$. Thus, we can define
\begin{align}
    \Prob_{\ast}(S) = \inf \left\{ \sum_{k=1}^{\infty}\abs{I_{k}} \, \Big| \, S \subseteq \bigcup_{k=1}^{\infty}I_{k}, \; I_{k} \subseteq [0,1] \right\}.
\end{align}
We have not verified that this is a valid definition of a probability measure, hence the subscript $\ast$. Via this, let us try computing $\Prob_{\ast}(\Q \cap [0,1])$. Enumerate the rationals in $[0,1]$ as $\{q_{1},q_{2},\ldots\}$. For any $\epsilon > 0$, we can choose intervals $I_{k} = [q_{k}-\epsilon/2^{k+1},q_{k}+\epsilon/2^{k+1}] \cap [0,1]$ for each $k \in \N$. Note that $\bigcup_{k=1}^{\infty}I_{k}$ contains all the rationals in $[0,1]$, and $\sum_{k=1} \abs{I_{k}} \leq \sum_{k=1}^{\infty}\epsilon/2^{k} = \epsilon$. Since $\epsilon$ was arbitrary, we have $\Prob_{\ast}(\Q \cap [0,1]) = 0$. This is consistent with our intuition that the probability of choosing a rational number from $[0,1]$ is $0$. So, $\Prob_{\ast}$ seems to be a reasonable definition of a probability measure on $([0,1],2^{[0,1]})$. However, we will show later that there exists a set $\cA \subseteq [0,1]$ such that $\Prob_{\ast}(\cA) = \Prob_{\ast}([0,1]\setminus \cA) = 1$, contradicting the property that $\Prob_{\ast}([0,1]) = 1$. We state the following theorem without proof.
\begin{theorem}
    There is no probability measure $\Prob$ on $([0,1],2^{[0,1]})$ such that $\Prob([a,b]) = b-a$ for all $[a,b] \subseteq [0,1]$.
\end{theorem}
If we remove the requirement that the probability measure of a subinterval $[a,b]$ is equal to its length, then we can define a probability measure on $([0,1],2^{[0,1]})$. For example, we can define $\Prob(S) = 1$ if $1/2 \in S$ and $\Prob(S) = 0$ otherwise. This is a valid probability measure, but it is not very useful. Keeping the requirement that the probability measure of a subinterval $[a,b]$ is equal to its length, we can instead restrict the collection of subsets of $[0,1]$ on which we define the probability measure. 

\begin{definition}
    A \eax{probability space} is a tuple $(\Omega,\cF,\Prob)$ where
    \begin{itemize}
        \item $\Omega$ is a non-empty set called the sample space,
        \item $\cF$ is a collection of subsets of $\Omega$ satisfying the following properties:
        \begin{itemize}
            \item $\Omega \in \cF$,
            \item if $A \in \cF$, then $\Omega \setminus A \in \cF$,
            \item if $A_{1},A_{2},\ldots \in \cF$, then $\bigcup_{i=1}^{\infty}A_{i} \in \cF$,
        \end{itemize}
        \item $\Prob:\cF\to[0,1]$ is a function satisfying:
        \begin{itemize}
            \item $\Prob(\Omega) = 1$,
            \item if $A_{1},A_{2},\ldots \in \cF$ are pairwise disjoint, then $\Prob\left(\bigcup_{i=1}^{\infty}A_{i}\right) = \sum_{i=1}^{\infty}\Prob(A_{i})$.
        \end{itemize}
    \end{itemize}
\end{definition}
Such an $\cF$ is termd a \eax{$\sigma$-algebra} on $\Omega$, and such a $\Prob$ is termed a \eax{probability measure} on $(\Omega,\cF)$.

\section{$\sigma$-algebras}

\begin{example}
    Given $\Omega$, trivial possible choices for $\cF$ are $\{\emptyset,\Omega\}$ and $2^{\Omega}$. A more interesting example is the following: suppose $\Omega$ is uncountable. Then $\cF = \{A \subseteq \Omega \mid A \text{ is countable or } \Omega \setminus A \text{ is countable}\}$ is a $\sigma$-algebra on $\Omega$.
\end{example}

Let $S \subseteq 2^{\Omega}$ be a collection of subsets of $\Omega$. We claim that there exists a notion of the smallest $\sigma$-algebra containing $S$. This is called the $\sigma$-algebra generated by $S$, and is denoted by $\sigma(S)$.

\begin{lemma}
    For any collection $S \subseteq 2^{\Omega}$, there exists a unique smallest $\sigma$-algebra containing $S$ satisfying
    \begin{align}
        \sigma(S) = \bigcap \{\cF \subseteq 2^{\Omega} \mid \cF \text{ is a $\sigma$-algebra and } S \subseteq \cF\}.
    \end{align}
\end{lemma}
\begin{proof}
    It is easy to check that $\sigma(S)$ is indeed a $\sigma$-algebra containing $S$. If $\cF$ is any $\sigma$-algebra containing $S$, then $\sigma(S) \subseteq \cF$. Thus, $\sigma(S)$ is the unique smallest $\sigma$-algebra containing $S$.
\end{proof}

\begin{definition}
    Let $(X,d)$ be a metric space (or a topological space), and let $\cS \subseteq 2^{X}$ be the colletion of all open sets in $X$. Then, the $\sigma$-algebra $\sigma(\cS)$ is called the \eax{Borel $\sigma$-algebra} on $X$, and its elements are called \eax{Borel set}s.
\end{definition}
We will show that there exists a probability measure on $([0,1],\cB)$, where $\cB$ is the Borel $\sigma$-algebra on $[0,1]$, which extends the notion of length of intervals to Borel sets.

Let us return to the problem of indefinite random walks on the line. Let $\Omega = \{\omega = (\omega_{1},\omega_{2},\ldots) \mid \omega_{i} \in \{0,1\}\}$, where $\omega_{i} = 0$ corresponds to a step to the left and $\omega_{i} = 1$ corresponds to a step to the right. If we let $A = \{0,1\}$, then $\Omega = A^{\N}$; consider the discrete topology on $A$, and the product topology on $\Omega$. If we let
\begin{align}
    \cS = \{U_{1} \times U_{2} \times \cdots \times U_{k} \times A \times A \cdots \mid U_{i} = \{0\} \text{ or } \{1\}\}
\end{align}
then one can show that the $\sigma$-algebra generated by $\cS$ is the Borel $\sigma$-algebra on $\Omega$. We can set $\cF = \sigma(\cS)$, and define a probability measure $\Prob$ on $(\Omega,\cF)$ by setting $\Prob(U_{1} \times U_{2} \times \cdots \times U_{k} \times A \times A \cdots) = 2^{-k}$ for all $U_{i} = \{0\}$ or $\{1\}$.
\\

\textit{July 24th.}
\begin{example}
    Let $(X,d)$ be a separable metric space (that is, there exists a countable dense subset). The Borel $\sigma$-algebra on $X$ is then
    \begin{align}
        \cB(X) = \sigma(\{B(p,r)\}_{p \in X, r >0}).
    \end{align}
    Note that the usual containment is $\sigma(\{B(p,r)\}_{p \in X, r >0}) \subseteq \cB(X)$ for a metric space $X$. So, it suffices to show that if $V \in \tau$, the topology on $X$, then $V \in \sigma(\{B(p,r)\}_{p \in X, r >0})$. Let $Q$ be a countable dense subset of $X$. Then, for any $V \in \tau$, we have
    \begin{align}
        V = \bigcup_{n=1}^{\infty} B(q_{n},r_{n}),
    \end{align}
    where $q_{n} \in Q \cap V$ and $r_{n} > 0$ (where $r_{n}$ is the largest possible radius such that $B(q_{n},r_{n}) \subseteq V$). Since $\sigma(\{B(p,r)\}_{p \in X, r >0})$ is a $\sigma$-algebra, we have $V \in \sigma(\{B(p,r)\}_{p \in X, r >0})$. Thus, we have $\cB(X) = \sigma(\{B(p,r)\}_{p \in X, r >0})$.
\end{example}

\begin{example}
    Let $(\Sigma,\cF,\Prob)$ be a probability space.
    
    \begin{itemize}
        \item $\cF$ is closed under countable unions, countable intersections, differences, and symmetric differences. Countable unions and intersections follow from the definition of a $\sigma$-algebra. If $A,B \in \cF$, then $A \setminus B = A \cap (\Omega \setminus B) \in \cF$. Finally, $A \triangle B = (A \setminus B) \cup (B \setminus A) \in \cF$.

        \item Let $(\Sigma,\cF,\Prob)$ be a probability space, and let $A_{1},A_{2},\ldots \in \cF$. Define
        \begin{align}
            \limsup_{n \to \infty} A_{n} \defeq \{x \in \Omega \mid x \text{ belongs to infinitely many } A_{n}\}
        \end{align}
        and
        \begin{align}
            \liminf_{n \to \infty} A_{n} \defeq \{x \in \Omega \mid x \text{ belongs to all but finitely many } A_{n}\}.
        \end{align}
        Then, both the $\limsup$ and $\liminf$ belong in $\cF$. To see this, simply note that
        \begin{align}
            \limsup_{n \to \infty} A_{n} = \bigcap_{n=1}^{\infty}\bigcup_{k=n}^{\infty}A_{k}, \quad \liminf_{n \to \infty} A_{n} = \bigcup_{n=1}^{\infty}\bigcap_{k=n}^{\infty}A_{k}.
        \end{align}

        \item One may show that $\Prob(\emptyset) = 0$, $\Prob(A) \leq \Prob(B)$ if $A \subseteq B$, $\Prob(A \cup B) = \Prob(A) + \Prob(B) - \Prob(A \cap B)$, and $\Prob(\bigcup A_{n}) \leq \sum \Prob(A_{n})$ for any countable collection of sets $A_{1},A_{2},\ldots \in \cF$.
        
        \item For a nested sequence of sets $A_{1} \subseteq A_{2} \subseteq \cdots$ in $\cF$, let $A = \bigcup_{n=1}^{\infty}A_{n}$ (this is usually denoted by $A_{n} \uparrow A$). Then, $\Prob(A) = \lim_{n \to \infty}\Prob(A_{n})$. To this end, let $B_{n} = A_{n}\setminus A_{n-1}$ for $n > 1$, and $B_{1} = A_{1}$. Then $A = \bigcup_{n=1}^{\infty}B_{n}$. Thus,
        \begin{align}
            \Prob(A) = \sum_{n=1}^{\infty}\Prob(B_{n}) = \lim_{N \to \infty}\sum_{n=1}^{N}\Prob(B_{n}) = \lim_{N \to \infty}\sum_{n=2}^{N} \Prob(A_{n}) - \Prob(A_{n-1}) + \Prob(A_{1}) = \lim_{N \to \infty}\Prob(A_{N}).
        \end{align}
    \end{itemize}    
\end{example}

\subsection{Other Systems of Sets}

The following are other interesting collection of subsets of $\Omega$ besides the Borel $\sigma$-algebra. 

\begin{enumerate}
    \item \eax{$\pi$-system}: a collection of subsets $S \subseteq \Omega$ such that $A,B \in S$ implies $A \cap B \in S$.
    \item \eax{$\lambda$-system}: a collection $\Lambda$ of subsets of $\Omega$ such that $\Omega \in \Lambda$, if $A,B \in \Lambda$ with $A \subseteq B$, then $B \setminus A \in \Lambda$, and if $A_{1},A_{2},\ldots \in \Lambda$, then $\bigcup_{n=1}^{\infty}A_{n} \in \Lambda$.
    \item \eax{algebra}: a collection $\cA$ of subsets of $\Omega$ such that $\Omega \in \cA$, if $A \in \cA$, then $\Omega \setminus A \in \cA$, and if $A,B \in \cA$, then $A \cup B \in \cA$.
\end{enumerate}

We had defined $\omega(S)$, the $\sigma$-algebra generated by a collection of subsets $S \subseteq 2^{\Omega}$. We can similarly define $\pi(S)$, the $\pi$-system generated by $S$, $\lambda(S)$, the $\lambda$-system generated by $S$, and $\cA(S)$, the algebra generated by $S$. 

\begin{theorem}[\eax{Sierpiński-Dynkin $\pi$-$\lambda$ Theorem}]
    Let $S,\cF \subseteq 2^{\Omega}$ be collections of subsets of $\Omega$. Then,
    \begin{enumerate}
        \item $\cF$ is a $\sigma$-algebra if and only if $\cF$ is both a $\pi$-system and a $\lambda$-system.
        \item If $S$ is a $\pi$-system, then $\lambda(S) = \sigma(S)$.
    \end{enumerate}
\end{theorem}
\begin{proof}
    \begin{enumerate}
        \item If $\cF$ is a $\sigma$-algebra, then it is trivially both a $\pi$-system and a $\lambda$-system. Conversely, if $\cF$ is both a $\pi$-system and a $\lambda$-system, we need to verify the axioms of a $\sigma$-algebra. We have $\Omega \in \cF$ since $\cF$ is a $\lambda$-system. If $A \in \cF$, then $\Omega \setminus A \in \cF$ since $\cF$ is a $\lambda$-system. Finally, if $A_{1},A_{2},\ldots \in \cF$, then we can write $B_{n} = \bigcup_{i=1}^{n} A_{i} \in \cF$ since $A_{i}^{c} \in \cF$ and $B_{n}^{c} = \bigcap_{i=1}^{n} A_{i}^{c} \in \cF$ since $\cF$ is a $\pi$-system. Then, $B_{n} \uparrow B = \bigcup_{i=1}^{\infty}A_{i}$, and since $\cF$ is a $\lambda$-system, we have $B \in \cF$. Thus, $\cF$ is a $\sigma$-algebra.
        \item Let $S$ be a $\pi$-system. Since $\lambda(S) \subseteq \sigma(S)$, it suffices to show that $\lambda(S)$ is a $\pi$-system. Fix $A \in S$, and define $\cC_{A} = \{B \in \lambda(S) \mid A \cap B \in \lambda(S)\}$. Note that $S \subseteq \cC_{A}$ because $S$ is a $\pi$-system. We show that $\cC_{A}$ is a $\lambda$-system. We have $\Omega \in \cC_{A}$ since $A \cap \Omega = A \in \lambda(S)$. If $B_{1},B_{2} \in \cC_{A}$ with $B_{1} \subseteq B_{2}$, then $A \cap B_{1} \subseteq A \cap B_{2}$, and since $A \cap B_{2} \in \lambda(S)$, we have $(A \cap B_{2})\setminus(A \cap B_{1}) = A \cap (B_{2}\setminus B_{1}) \in \lambda(S)$, and thus $B_{2}\setminus B_{1} \in \cC_{A}$. Finally, if $B_{1},B_{2},\ldots$ are in $\cC_{A}$, then $A \cap B_{n} \in \lambda(S)$ for all $n$, and thus $\bigcup (A\cap B_{n}) = A\cap (\bigcup B_{n}) \in \lambda(S)$, and hence $\bigcup B_{n} \in \cC_{A}$. Thus, $\cC_{A}$ is a $\lambda$-system. Since $\lambda(S)$ is the smallest $\lambda$-system containing $S$, we have $\lambda(S) = \cC_{A}$, and hence for any $B,C \in S$, we have $B,C \in S$ implies $B\cap C = A\cap C \in S$. Thus, $\lambda(S)$ is a $\pi$-system.
    \end{enumerate}
\end{proof}


\begin{lemma}
    Let $S$ be a $\pi$-system, and let $\Prob$ and $\Q$ be two probability measures on $(\Sigma,\sigma(S))$ such that $\Prob(A) = \Q(A)$ for all $A \in S$. Then, $\Prob(A) = \Q(A)$ for all $A \in \sigma(S)$.
\end{lemma}
\begin{proof}
    Let $\cC = \{B \in \sigma(S) \mid \Prob(B) = \Q(B)\}$. Then, $S \subseteq \cC$, and we show that $\cC$ is a $\lambda$-system. We have $\Omega \in \cC$ since $\Prob(\Omega) = 1 = \Q(\Omega)$. If $B_{1},B_{2} \in \cC$ with $B_{1} \subseteq B_{2}$, then $\Prob(B_{2}\setminus B_{1}) = \Prob(B_{2}) - \Prob(B_{1}) = \Q(B_{2}) - \Q(B_{1}) = \Q(B_{2}\setminus B_{1})$, and hence $B_{2}\setminus B_{1} \in \cC$. Finally, if $B_{1},B_{2},\ldots$ are in $\cC$, then $\Prob(A) = \lim_{n \to \infty} \Prob(A_n)$, and since $\Q(A) = \lim_{n \to \infty} \Q(A_n)$, we have $\Prob(A) = \Q(A)$.
\end{proof}

\textit{July 28th.}

\begin{example}
    We wish to obtain an example of $S$ which is \textit{not} a $\pi$-system, and two probability measures $\Prob$ and $\Q$ on $(\Sigma,\sigma(S))$ such that $\Prob(A) = \Q(A)$ for all $A \in S$, but $\Prob(A) \neq \Q(A)$ for some $A \in \sigma(S)$. Let $\Omega = \{1,2,3,4\}$, and let $S = \{\{1,2\},\{2,3\}\}$. Then, $\sigma(S) = 2^{\Omega}$. Define $\Prob$ and $\Q$ on $(\Omega,\sigma(S))$ by
    \begin{align}
        \Prob(\{1\}) = \Prob(\{2\}) = \Prob(\{3\}) = \Prob(\{4\}) = 1/4,
    \end{align}
    and
    \begin{align}
        \Q(\{1\}) = \Q(\{3\}) = 1/6, \quad \Q(\{2\}) = \Q(\{4\}) = 1/3.
    \end{align}
    Then, $\Prob(\{1,2\}) = \Q(\{1,2\}) = 1/2$ and $\Prob(\{2,3\}) = \Q(\{2,3\}) = 1/2$, but $\Prob(\{3\}) = 1/4 \neq 1/6 = \Q(\{3\})$.
\end{example}

\begin{example}
    Let $\Omega \neq \emptyset$, and $A_{\lambda}$ be a collection of subsets of $\Omega$ indexed by $\lambda \in \Lambda$. Consider the $\sigma$-algebra $\sigma(\{A_{\lambda}\}_{\lambda \in \Lambda})$, and pick $B \in \sigma(\{A_{\lambda}\}_{\lambda \in \Lambda})$. We will show that there exists a countable subcollection $\{A_{\lambda_{n}}\}_{n=1}^{\infty}$ such that $B \in \sigma(\{A_{\lambda_{n}}\}_{n=1}^{\infty})$. Let $\cC = \{B \in \sigma(\{A_{\lambda}\}_{\lambda \in \Lambda}) \mid B \in \sigma(\{A_{\lambda_{n}}\}_{n=1}^{\infty}) \text{ for some countable subcollection } \{A_{\lambda_{n}}\}_{n=1}^{\infty}\}$. Then, $\cC$ is a $\sigma$-algebra containing $\{A_{\lambda}\}_{\lambda \in \Lambda}$, and hence $\sigma(\{A_{\lambda}\}_{\lambda \in \Lambda}) \subseteq \cC$. Thus, $B \in \cC$, and hence there exists a countable subcollection $\{A_{\lambda_{n}}\}_{n=1}^{\infty}$ such that $B \in \sigma(\{A_{\lambda_{n}}\}_{n=1}^{\infty})$.
\end{example}

\section{The Lebesgue Measure}

\begin{theorem}
    There exists a unique Borel measure $\lambda$ on $[0,1]$ such that $\lambda(I) = \abs{I}$ for all intervals $I \subseteq [0,1]$.
\end{theorem}

\begin{proof}
    For the $\pi$-system $S = \{(a,b] \mid 0 \leq a < b \leq 1\}$, we know that $\sigma(S)$ is a Boreal $\sigma$-algebra on $[0,1]$. If $\lambda$ and $\mu$ are two Borel measures on $[0,1]$ such that $\lambda(I) = \mu(I) = \abs{I}$ for all intervals $I \subseteq [0,1]$, then by the previous lemma, we have $\lambda(A) = \mu(A)$ for all $A \in \sigma(S)$. Thus, the Borel measure is unique.

    To show the existence of such a measure, we will define a pseudomeasure $\lambda_{\ast}$ that satisfies the interval property, and then show that it is a measure on the Borel $\sigma$-algebra. For any $E \subseteq [0,1]$, define
    \begin{align}
        \lambda_{\ast}(E) = \inf\left\{ \sum_{i=1}^{\infty} \abs{I_{i}} \, \Big| \, I_{i} \text{ are intervals and } E \subseteq \bigcup_{i=1}^{\infty} I_{i} \right\}
    \end{align}
    Infer that $0 \leq \lambda_{\ast}(E) \leq 1$ for all $E \subseteq [0,1]$. Moreover, let $A$ and $B$ be two sets in $[0,1]$. Then, for any $\eps > 0$, we have
    \begin{align}
        \lambda_{\ast}(A \cup B) \leq \sum_{k=1}^{\infty} \abs{I_{k}} + \sum_{k=1}^{\infty} \abs{J_{k}} \leq \lambda_{\ast}(A) + \eps + \lambda_{\ast}(B) + \eps \implies \lambda_{\ast}(A \cup B) \leq \lambda_{\ast}(A) + \lambda_{\ast}(B).
    \end{align}
    Finally, let $(a,b] \subseteq [0,1]$ be an interval. Then, $\lambda_{\ast}((a,b]) = b-a$. To see this, suppose $(a,b] \subseteq \bigcup_{k=1}^{\infty} I_{k}$ for some intervals $I_{k}$. Now consider $[a+\eps,b]$ for some $\eps > 0$ such that $[a+\eps,b] \subseteq \bigcup_{k=1}^{\infty} I_{k}$. This is a compact subinterval, so we can extract a finite subcover $[a+\eps,b] \subseteq \bigcup_{k=1}^{N} I_{k}$. Then,
    \begin{align}
        b-a-\eps = \abs{[a+\eps,b]} \leq \sum_{k=1}^{N} \abs{I_{k}} \leq \sum_{k=1}^{\infty} \abs{I_{k}} \implies b-a \leq \lambda_{\ast}((a,b]).
    \end{align}
    Now, since $(a,b] \subseteq (a,b+\eps)$, a similar line of reasoning shows that $\lambda_{\ast}((a,b]) \leq b-a$. Thus, we have $\lambda_{\ast}((a,b]) = b-a$. The property of finite subadditivity also implies the property of countable subadditivity (show this). We now construct the $\sigma$-algebra $\cF$ on which $\lambda_{\ast}$ is a measure. We pick
    \begin{align}
        \cF = \{A \subseteq [0,1] \mid \lambda_{\ast}(E) = \lambda_{\ast}(E \cap A) + \lambda_{\ast}(E \cap A^{c}) \text{ for all } E \subseteq [0,1]\}.
    \end{align}
    We show that $\cF$ is a $\sigma$-algebra containing the Borel $\sigma$-algebra, and that $\lambda_{\ast}$ is a measure on $\cF$. It is left as a minor exercise to show that $\cF$ is closed under intersections and complements (making $\cF$ a $\pi$-system). We show $\cF$ is a $\lambda$-system. Let $A,B \in \cF$. Then $B \setminus A = B \cap A^{c}$ is in $\cF$. Then, for any $E \subseteq [0,1]$, we have
    \begin{align}
        \lambda_{\ast}(E)
        &= \lambda_{\ast}(E\cap A_n) + \lambda_{\ast}(E\cap A_n^{c}) \\
        &\ge \lambda_{\ast}(E\cap A_n) + \lambda_{\ast}(E\cap A^{c}),
    \end{align}
    because $A^{c} \subseteq A_n^{c}$. Since $A_n \uparrow A$, we have $E\cap A_n \uparrow E\cap A$, and hence
    \[
    \lambda_{\ast}(E\cap A_n) \uparrow \lambda_{\ast}(E\cap A).
    \]
    Passing to the limit gives
    \[
    \lambda_{\ast}(E) \ge \lambda_{\ast}(E\cap A) + \lambda_{\ast}(E\cap A^{c}).
    \]
    On the other hand, subadditivity of the outer measure gives
    \[
    \lambda_{\ast}(E) \le \lambda_{\ast}(E\cap A) + \lambda_{\ast}(E\cap A^{c}).
    \]
    Therefore,
    \[
    \lambda_{\ast}(E)=\lambda_{\ast}(E\cap A)+\lambda_{\ast}(E\cap A^{c}),
    \]
    so $A\in\cF$.
\end{proof}

\noindent \textit{July 30th.}

Note that $\cF \subsetneq 2^{[0,1]}$; that is, there exist subsets of $[0,1]$ that are not in $\cF$. We shall assume that $\lambda_{\ast}(E) = \lambda_{\ast}(E+x)$ for all $E \subseteq [0,1]$ and $x \in \R$ such that $E+x \subseteq [0,1]$. This is a reasonable assumption, since we expect the Lebesgue measure to be translation invariant (one can prove this). Moreover, instead of looking at $[0,1]$ we shall look at $([0,1),\oplus)$, where $\oplus$ is addition modulo $1$. Now suppose we manage to break $[0,1)$ into countably many disjoint pieces $A_{1},A_{2},\ldots$ such that $\bigcup_{n=1}^{\infty}A_{n} = [0,1)$, and $A_{i} = A_{1} \oplus r_{i}$ for some $r_{i} \in [0,1)$. Then, we would have $\lambda_{\ast}(A_{i}) = \lambda_{\ast}(A_{1})$ for all $i$, and hence $\lambda_{\ast}([0,1)) = \sum_{n=1}^{\infty}\lambda_{\ast}(A_{n}) = \sum_{n=1}^{\infty}\lambda_{\ast}(A_{1})$. This is a contradiction, since $\lambda_{\ast}([0,1)) = 1$ and $\sum_{n=1}^{\infty}\lambda_{\ast}(A_{1})$ is either $0$ or $\infty$. Thus, we have shown that there exists a set $A \subseteq [0,1)$ such that $A \notin \cF$.

To construct these $A_{i}$'s, simply look at the equivalence relation $\sim$ on $[0,1)$ defined by $x \sim y$ if and only if $x-y \in \Q$. Then, we can pick a representative from each equivalence class to form a set $A_{1}$. Then, we can define $A_{i} = A_{1} \oplus r_{i}$ for some $r_{i} \in [0,1)$ such that $\{r_{i}\}_{i=1}^{\infty}$ is a complete set of representatives of $\Q \cap [0,1)$. This gives us the desired decomposition of $[0,1)$ into disjoint pieces.

\subsection{Carathéodory's Extension Theorem}
\textit{August 4th.}

Suppose $\Omega$ is non-empty, and $\cA$ is an \eax{algebra}; that is, $\cA$ is a collection of subsets of $\Omega$ such that $\Omega \in \cA$, and $\cA$ is closed under complements and finite unions. Further say that $\mu$ is a \eax{pre-measure} on $\cA$; that is, $\mu:\cA \to [0,1]$ satisfies $\mu(\emptyset) = 0$, and if $A_{1},A_{2},\ldots \in \cA$ are pairwise disjoint, then $\mu(\bigcup_{i=1}^{\infty}A_{i}) = \sum_{i=1}^{\infty}\mu(A_{i})$. Then, we can extend $\mu$ to a measure on the $\sigma$-algebra generated by $\cA$. This is the content of \eax{Carathéodory's extension theorem}.

\begin{remark}
    \begin{enumerate}
        \item Note that since $\cA$ is a $\pi$-system, such an extension must be unique.
        \item For a $\pi$-system $S$, a pre-measure on $\S$ gives a pre-measure on $\cA(S)$, the algebra generated by $S$, since every set $B \in \cA(S)$ can be written as a finite disjoint union of sets in $S$.
    \end{enumerate}
\end{remark}

We now wish to assign a `correct' setting to talk about infinite coin tosses (or simple random walks on $\Z$). Let $\Omega = \{0,1\}^{\N}$, and we define
\begin{align}
    S = \{\{\bv{x} \mid x_{1} = a_{1},\, x_{2} = a_{2},\ldots,x_{n} = a_{n} \} \mid n \in \N, \, a_{i} \in \{0,1\}\} \cup \{\emptyset\}
\end{align}
and let $\cF = \sigma(S)$. Verify that $S$ is a $\pi$-system; we can set
\begin{align}
    \mu(\{\bv{x} \mid x_{1} = a_{1},\, x_{2} = a_{2},\ldots,x_{n} = a_{n} \}) = \prod_{i=1}^{n} p^{a_{i}}(1-p)^{1-a_{i}}.
\end{align}
For now, assume that $\mu$ is a pre-measure on $S$. Then there'll exist a unique probability measure on $(\Omega,\sigma(S))$ that extends $\mu$.

\section{Measurable Maps and Random Variables}

Note that a space $(\Omega,\cF)$ is called a \eax{measurable space}.
\begin{definition}
    A map $T:(\Omega,\cF) \to (\tilde{\Omega},\tilde{\cF})$ between two measurable spaces is termed a \eax{measurable function} if $T^{-1}(A) \in \cF$ for all $A \in \tilde{\cF}$. When $(\tilde{\Omega},\tilde{\cF}) = (\R,\cB)$, we call $T$ a \eax{random variable}.
\end{definition}
In most situations, $X$ is usually a topological space and $\cB(X)$ is the Borel $\sigma$-algebra on $X$. For the case of $X = \R$, a random variable can be interpreted as a meaningful quantity which can be studied. 

\begin{proposition}
    $T:(\Omega,\cF) \to (\R,\cB)$ is measurable if and only if $T^{-1}(U) \in \cF$ whenever $U$ is open in $\R$.
\end{proposition}
\begin{proof}
    One direction is clear; for the converse direction, let $\tilde{\cF} = \{W \subseteq \R \mid T^{-1}(W) \in \cF\}$. It is relatively easy to show that $\tilde{\cF}$ is a $\sigma$-algebra on $\R$. Thus, since $\cB \subseteq \tilde{\cF}$, we have $T^{-1}(A) \in \cF$ for all $A \in \cB$.
\end{proof}

Note that the image space being in $\R$ did not play a role in the above proof; for any image space $\tilde{\Omega}$, we can define $\tilde{\cF} = \{W \subseteq \tilde{\Omega} \mid T^{-1}(W) \in \cF\}$, and show that $\tilde{\cF}$ is a $\sigma$-algebra on $\tilde{\Omega}$. Such a $\tilde{\cF}$ is called the \eax{pushforward $\sigma$-algebra} of $\cF$ under $T$, and is the largest $\sigma$-algebra on $\tilde{\Omega}$ such that $T$ is measurable. One can think in the converse way, where we are given a map $T:\Omega \to (\tilde{\Omega},\tilde{F})$; to make such a $T$ measurable, one can simply equip $\Omega$ with the powerset $\sigma$-algebra $2^{\Omega}$, and then $T$ is automatically measurable. Alternatively, we can define $\cF = \{T^{-1}(A) \mid A \in \tilde{\cF}\}$, and then $T$ is measurable. This is called the \eax{pullback $\sigma$-algebra} of $\tilde{\cF}$ under $T$, and is the smallest $\sigma$-algebra on $\Omega$ such that $T$ is measurable.

The most non-trivial yet basic example of a random variable is the \eax{indicator function} of a set $A \in \cF$, which is defined as $\1_{A}(\omega) = 1$ if $\omega \in A$, and $0$ otherwise. It is easy to check that $\1_{A}$ is a random variable.

\begin{lemma}
    If $X,Y:(\Omega,\cF) \to (\R,\cB)$ are random variables, then $X+Y$ is also a random variable. More generally, if $f:\R^{2} \to \R$ is continuous, then $f(X,Y)$ is a random variable.
\end{lemma}

\textit{August 7th.}

\begin{example}
    Consider the infinite simple random walk from before; for $a \in \Z$, define $\tau_{a} = \inf\{n > 0 \mid S_{n}=a\}$, the first hitting time of $a$. Then we claim that $\tau_{a}$ is measurable and, hence, a random variable. To see this, note that $\tau_{a}^{-1}((n,\infty)) = \{\omega \in \Omega \mid \tau_{a}(\omega) > n\} = \{\omega \in \Omega \mid S_{1}(\omega) \neq a, S_{2}(\omega) \neq a, \ldots, S_{n}(\omega) \neq a\} = \bigcap_{k=1}^{n}\{\omega \in \Omega \mid S_{k}(\omega) \neq a\}$, which is in $\cF$ since $S_{k}$ is measurable for all $k$. Thus, $\tau_{a}$ is measurable.
\end{example}

\begin{definition}
    Let $(\Omega,\cF,\Prob)$ be a probability space, and let $T:(\Omega,\cF) \to (\tilde{\Omega},\tilde{\cF})$ be a measurable map. Then, the \eax{pushforward measure} of $\Prob$ under $T$ is the measure $\tilde{\Prob}$ on $(\tilde{\Omega},\tilde{\cF})$ defined by $\tilde{\Prob}(A) = \Prob(T^{-1}(A))$ for all $A \in \tilde{\cF}$.
\end{definition}
Indeed, one can show that the pushforward measure $\tilde{\Prob}$ is a probability measure on $(\tilde{\Omega},\tilde{\cF})$.

\begin{definition}
    Given a Borel probability measure $\mu$ on $(\R,\cB)$, we define the \eax{cumulative distribution function} of $\mu$ to be the function $F_{\mu}:\R \to [0,1]$ defined by $F_{\mu}(x) = \mu((-\infty,x])$ for all $x \in \R$. 
\end{definition}

\begin{proposition}
    With $\mu$ as above, the following hold true.
    \begin{enumerate}
        \item $F_{\mu}$ is non-decreasing.
        \item $\lim_{x \to \infty} F_{\mu}(x) = 1$ and $\lim_{x \to -\infty} F_{\mu}(x) = 0$.
        \item $F_{\mu}$ is right-continuous.
    \end{enumerate}
\end{proposition}

\begin{proof}
    \begin{enumerate}
        \item If $x_{1} \leq x_{2}$, then $(-\infty,x_{1}] \subseteq (-\infty,x_{2}]$, and hence $F_{\mu}(x_{1}) = \mu((-\infty,x_{1}]) \leq \mu((-\infty,x_{2}]) = F_{\mu}(x_{2})$.
        \item Note that $\R = \bigcup_{n=1}^{\infty}(-\infty,n]$; defining $A_{k} = (-\infty,k]$, we have $A_{1} \subseteq A_{2} \subseteq \cdots$ and $A_{k} \uparrow \R$. Since $\mu$ is a measure, $\mu(A_{k}) \uparrow \mu(\R) = 1$, and hence $\lim_{x \to \infty} F_{\mu}(x) = 1$. Similarly, $\lim_{x \to -\infty} F_{\mu}(x) = 0$.
        \item Let $x_{n}$ be a decreasing sequence of numbers such that $x_{n} \downarrow x$. Then, $(-\infty,x_{1}] \supseteq (-\infty,x_{2}] \supseteq \cdots$ and $(-\infty,x_{n}] \downarrow (-\infty,x]$. Since $\mu$ is a measure, $\mu((-\infty,x_{n}]) \downarrow \mu((-\infty,x])$, and hence $F_{\mu}(x_{n}) \downarrow F_{\mu}(x)$. Thus, $F_{\mu}$ is right-continuous.
    \end{enumerate}
\end{proof}

\begin{theorem}
    Suppose $F:\R \to [0,1]$ is a non-decreasing, right-continuous function such that $\lim_{x \to -\infty} F(x) = 0$ and $\lim_{x \to \infty} F(x) = 1$. Then, there exists a unique Borel probability measure $\mu$ on $(\R,\cB)$ such that $F = F_{\mu}$.
\end{theorem}
Thus, there is a one-to-one correspondence between Borel probability measures on $(\R,\cB)$ and cumulative distribution functions. In general, it easier to study the cumulative distribution function of a probability measure than the measure itself.

\subsection{The Pushfoward Measure}
\textit{August 11th.}

In many situations, we can model a random experiment in several ways. Let $\Omega = [0,1]$, $\cF = \cB$, and $\lambda$ the Lebesgue measure. Define the random variables $X_{1} = \1_{[0,1/2]}$ and $X_{2} = \1_{[0,1/4] \cup [1/2,3/4]}$. Consider $\Gamma = (X_{1},X_{2}):\Omega \to \R^{2}$, and let $\mu = \lambda \circ \Gamma^{-1}$ be the pushforward measure of $\lambda$ under $\Gamma$. The support of $\mu$ is the set $\Gamma(\Omega) = \{(1,0),(0,1),(1,1),(0,0)\}$. The measure at each of these points is given by $\mu(\{(1,0)\}) = \lambda(\Gamma^{-1}(\{(1,0)\})) = \lambda([0,1/4]) = 1/4$, $\mu(\{(0,1)\}) = \lambda(\Gamma^{-1}(\{(0,1)\})) = \lambda([1/2,3/4]) = 1/4$, $\mu(\{(1,1)\}) = \lambda(\Gamma^{-1}(\{(1,1)\})) = \lambda([1/4,1/2]) = 1/4$, and $\mu(\{(0,0)\}) = \lambda(\Gamma^{-1}(\{(0,0)\})) = \lambda([3/4,1]) = 1/4$. Thus, we have a probability measure on the support of $\mu$. That is, we have modeled tossing a fair coin twice, without explicitly defining the sample space as $\{0,1\}^{2}$. The probability space does not matter in most contexts; the random variables and their distributions are what matter.

Naturally, we can ask further questions. Whether the following models can exist: tossing a coin $n$ times, and tossing a coin countably many times. The first question is easy to answer; we can define the $k$-th random variable $X_{k}$ exactly as
\begin{align}
    X_{k} = \1_{[0,\frac{1}{2^{k}}] \cup [\frac{2}{2^{k}},\frac{3}{2^{k}}] \cup \cdots \cup [\frac{2^{k}-2}{2^{k}},\frac{2^{k}-1}{2^{k}}]}.
\end{align}
We leave it as an exercise to verify that $(X_{1},X_{2},\ldots,X_{n})$ has the same distribution as tossing a fair coin $n$ times.

\begin{theorem}
    Given any probability measure $\mu$ on $(\R,\cB)$, there exists a measurable map $T:([0,1],\cB([0,1]),\lambda) \to (\R^{n},\cB(\R^{n}))$ such that $\mu = \lambda \circ T^{-1}$. That is, the pushforward measure of $\lambda$ under $T$ is $\mu$. 
\end{theorem}

Before we proceed with this theorem, we show the following proposition.
\begin{proposition}
    For any Borel probability measures $\mu$ and $\nu$ on $(\R,\cB)$, if $F_{\mu} = F_{\nu}$, then $\mu = \nu$. That is, the cumulative distribution function uniquely determines the probability measure.
\end{proposition}
\begin{proof}
    Let $\mu$ and $\nu$ be the required measures such that $F_{\mu} = F_{\nu}$, or $\mu((-\infty,x]) = \nu((-\infty,x])$ for all $x \in \R$. Let $S = \{(-\infty,x] \mid x \in \R\}$, which is a $\pi$-system. Note that $\sigma(S) = \cB$; thus, by the Sierpiński-Dynkin $\pi$-$\lambda$ theorem, we have $\mu(A) = \nu(A)$ for all $A \in \cB$. Hence, $\mu = \nu$.
\end{proof}

We now prove \textbf{Theorem 1.23} and \textbf{Theorem 1.24} together.
\begin{proof}[Proof of theorems]
    Suppose $F:\R \to [0,1]$ be a non-decreasing, right-continuous function such that $\lim_{x \to -\infty} F(x) = 0$ and $\lim_{x \to \infty} F(x) = 1$. We will show that there exists a measurable map $T:([0,1],\cB([0,1]),\lambda) \to (\R,\cB)$ such that $F_{\lambda \circ T^{-1}} = F$; note that this implies the second theorem by the previous proposition.

    For such an $F$, we can define a generalized inverse $T$ defined as $T(x) = \inf\{p \in \R \mid F(p) \geq x\}$. For any $x \in [0,1]$, the set $\{p \in \R \mid F(p) \geq x\}$ is non-empty (since $\lim_{p \to \infty} F(p) = 1$) and has infimum, so $T$ is well-defined. Moreover, $T$ is non-decreasing and takes values in $\R$. To show measurability of $T$, note that for any $a \in \R$, we have
    \begin{align}
        \{x \in [0,1] \mid T(x) > a\} = \{x \in [0,1] \mid F(a) < x\} = (F(a), 1] \cap [0,1],
    \end{align}
    which is a Borel set in $[0,1]$. Thus, $T$ is measurable. We now verify that $F_{\lambda \circ T^{-1}} = F$; For any $x \in \R$, we have
    \begin{align}
        (\lambda \circ T^{-1})((-\infty, x]) 
        &= \lambda(T^{-1}((-\infty, x])) 
        = \lambda(\{u \in [0,1] \mid T(u) \leq x\}) \notag \\
        &= \lambda(\{u \in [0,1] \mid \inf\{p \in \R \mid F(p) \geq u\} \leq x\})
        = \lambda(\{u \in [0,1] \mid F(x) \geq u\}) \notag \\
        &= \lambda([0, F(x)])
        = F(x).
    \end{align}
    Thus, $F_{\lambda \circ T^{-1}} = F$, establishing the existence of the desired measure.
\end{proof}

\section{Probability Measures on $\R$}
\textit{August 14th.}

We first wish to generalize the notion of the integral of a function with respect to the Lebesgue measure. A function $\varphi:[a,b] \to \R$ is said to be a \eax{simple function} if it is of the form $\varphi = \sum_{i=1}^{n} a_{i}\1_{A_{i}}$ for some $n \in \N$, $a_{i} \in \R$, and $A_{i} \in \cB$. That is, $\varphi$ takes finitely many values, and the pre-image of each value is a Borel set. For a simple function $\varphi$, we define the integral of $\varphi$ to be
\begin{align}
    \int_{[a,b]} \varphi(x) \, dx = \sum_{i=1}^{n} a_{i}\lambda(A_{i}).
\end{align}
In general, a function $g:[a,b] \to \R$ is said to be \eax{Borel measurable} if $g^{-1}(A) \in \cB$ for all Borel sets $A \subseteq \R$. Note that any simple function is Borel measurable. The integral of a non-negative Borel measurable function $g:[a,b] \to \R$ is then defined as
\begin{align}
    \int_{[a,b]} g(x) \, dx = \sup\left\{\int_{[a,b]} \varphi(x) \, dx \; \Big| \; 0 \leq \varphi \leq g, \, \varphi \text{ is simple}\right\}.
\end{align}
One can infer that $\int(f+\alpha g) = \int f + \alpha \int g$, $g \leq f \implies \int g \leq \int f$, and $\int_{I_{1} \sqcup I_{2}} g = \int_{I_{1}} g + \int_{I_{2}} g$ for disjoint intervals $I_{1}$ and $I_{2}$, and non-negative Borel measurable functions $f$ and $g$. We can extend the definition of the integral to all Borel measurable functions $g$ by defining $\int g = \int g^{+} - \int g^{-}$, where $g^{+} = \max\{g,0\}$ and $g^{-} = \max\{-g,0\}$. Thus, we are fit to define the following.

\begin{definition}
    A probability measure $\mu \in \cP(\R)$ is said to have a \eax{density} $f$ if $f:\R \to \R$ is a non-negative Borel measurable function such that $\int_{\R} f(x) \, dx = 1$ and $F_{\mu}(x) = \int_{-\infty}^{x} f(t) \, dt$ for all $x \in \R$. 
\end{definition}

Suppose $\mu \in \cP(\R)$ has a density $f$. Then, for any Borel set $A \subseteq \R$, we have
\begin{align}
    \mu(A) = \int_{A} f(x) \, dx.
\end{align}
Note that not every probability measure has a density; for example, the measure $\mu = \delta_{0}$, indicator of the set $\{0\}$, does not have a density. In fact, if there exists $p \in \R$ with $\mu(\{p\}) > 0$, then $\mu$ does not have a density. This is because if $\mu$ has a density $f$, then $\mu(\{p\}) = \int_{\{p\}} f(x) \, dx = 0$, which is a contradiction. Thus, any probability measure with a density must be non-atomic; that is, it cannot assign positive measure to any singleton set. However, the converse is not true; there exist non-atomic probability measures that do not have a density. For example, consider the Cantor measure $\mu$ on $[0,1]$, which is non-atomic, but does not have a density.

We now wish to define a metric on the space of probability measures on $\R$.
\begin{enumerate}
    \item For any $\mu,\nu \in \cP(\R)$, we define the \eax{total variation distance} between $\mu$ and $\nu$ to be
    \begin{align}
        D_{1}(\mu,\nu) = d_{TV}(\mu,\nu) \defeq \sup_{A \in \cB} \abs{\mu(A) - \nu(A)}.
    \end{align}
    \item The CDFs may be utilized to define
    \begin{align}
        D_{2}(\mu,\nu) \defeq \sup_{x \in \R} \abs{F_{\mu}(x) - F_{\nu}(x)}.
    \end{align}
\end{enumerate}

Both of these metrics are fine as they are, but fail in being `nice' with respect to convergence of measures; the sequence of measures $\{\mu_{n} = \delta_{1/n}\}_{n=1}^{\infty}$ fails to converge to $\delta_{0}$ in either of these metrics, even though we expect it to converge intuitively.
\\

\noindent \textit{August 18th.}

Hereforth, we define the \eax{Levy distance} on $\cP(\R)$ as
\begin{align}
    d(\mu,\nu) \defeq \inf\{\epsilon > 0 \mid F_{\mu}(t+\epsilon)+\epsilon \geq F_{\nu}(t), \; F_{\nu}(t+\epsilon)+\epsilon \geq F_{\mu}(t) \text{ for all } t \in \R\}.
\end{align}
In fact, this can be generalized to a metric on Borel probability measures on any metric space $(X,\cD)$ as follows (this is termed the \eax{Levy-Prokhorov metric}).
\begin{align}
    d_{LP}(\mu,\nu) \defeq \inf\{\epsilon > 0 \mid \mu(A+\epsilon)+\epsilon \geq \nu(A), \; \nu(A+\epsilon)+\epsilon \geq \mu(A) \text{ for all closed } A \subseteq X\}
\end{align}
where the \eax{$\epsilon$-fattening} of $A$ is defined as
\begin{align}
    A+\epsilon \defeq \{x \in X \mid \text{ there exists } y \in A \text{ such that } \cD(x,y) < \epsilon\}.
\end{align}
Note that these two metrics are not the same on $\cP(\R)$, but are equivalent. Also, via the Levy distance, we can show $d(\delta_{1/n},\delta_{0}) = 1/n$, showing that $\delta_{1/n} \to \delta_{0}$ in this metric.
\\

\textit{August 25th.}

\begin{theorem}
    Suppose $\mu,\mu_{1},\mu_{2},\ldots \in \cP(\R)$ are Borel probability measures. Then $d(\mu_{n},\mu) \xrightarrow{n \toinf} 0$ if and only if $F_{\mu_{n}}(x) \xrightarrow{n \toinf} F_{\mu}(x)$ for all continuity points $x \in \R$ of $F_{\mu}$.
\end{theorem}
\begin{proof}
    We first show the forward implication. Suppose $d(\mu_{n},\mu) \to 0$. Then, for any $\epsilon > 0$, there exists $N_{\epsilon} \in \N$ such that $d(\mu_{n},\mu) < \epsilon$ for all $n \geq N_{\epsilon}$. Now fix $\epsilon > 0$ and $x \in \R$. Then, for all $n \geq N_{\epsilon}$, we have $F_{\mu}(x+\epsilon) + \epsilon \geq F_{\mu_{n}}(x)$. We can replace the right hand side with the limit superior to get
    \begin{align}
        F_{\mu}(x+\epsilon) + \epsilon \geq \limsup_{n \toinf} F_{\mu_{n}}(x) \quad \text{ for all } \epsilon > 0.
    \end{align}
    This implies that $F_{\mu}(x) \geq \limsup_{n \toinf} F_{\mu_{n}}(x)$. Similarly, for all $n \geq N_{\epsilon}$, we have $F_{\mu_{n}}(x) + \epsilon \geq F_{\mu}(x-\epsilon)$. By a similar argument, we have $\liminf_{n \toinf} F_{\mu_{n}}(x) \geq F_{\mu}(x)$ if $x$ is a continuity point of $F_{\mu}$. Thus, we have shown that $\limsup_{n \toinf} F_{\mu_{n}}(x) \leq F_{\mu}(x) \leq \liminf_{n \toinf} F_{\mu_{n}}(x)$ or $F_{\mu_{n}}(x) \to F_{\mu}(x)$ for all continuity points.

    For the converse, if $\cC$ denotes the set of continuity points of $F_{\mu}$, then $\cC$ is cocountable and hence dense in $\R$. Now fix $\eps \in (0,1)$. Let $A,B \in \R$ be numbers such that $F_{\mu}(A) < \eps$ and $F_{\mu}(B) \geq 1-\eps$. Pick points $\{x_{i}\}_{i=0}^{m}$ in $[A,B] \cap \cC$ such that $A=x_{0} < x_{1} < \cdots < x_{m} = B$ and $x_{i+1}-x_{i} < \eps$ for all $i$. By the choice of the $x_{i}$'s, we have $F_{\mu_{n}}(x_{i}) \to F_{\mu}(x_{i})$ for all $i$, and so there exists $N_{\eps} \in \N$ such that $\abs{F_{\mu_{n}}(x_{i}) - F_{\mu}(x_{i})} < \eps$ for all $n \geq N_{\eps}$ and all $i$. Now fix $t \in [A,B]$; so $t \in [x_{i},x_{i+1}]$ for some $i$. Then, for all $n \geq N_{\eps}$, we have
    \begin{align}
        F_{\mu_{n}}(t+\eps) \geq F_{\mu_{n}}(x_{i+1}) \geq F_{\mu}(x_{i+1}) - \eps \geq F_{\mu}(t) - \eps \implies F_{\mu_{n}}(t+\eps) + \eps \geq F_{\mu}(t)
    \end{align}
    for all $n \geq N_{\eps}$ and $t \in [A,B]$. If $t > B$, then both $F_{\mu_{n}}(t)$ and $F_{\mu}(t)$ are at least $1-2\eps$, and hence $F_{\mu_{n}}(t+\eps) + \eps \geq F_{\mu}(t)$. If $t < A$, then both $F_{\mu_{n}}(t)$ and $F_{\mu}(t)$ are at most $2\eps$, and hence $F_{\mu_{n}}(t+\eps) + \eps \geq F_{\mu}(t)$. Thus, we have shown that for all $\eps > 0$, there exists $N_{\eps} \in \N$ such that for all $n \geq N_{\eps}$, we have
    \begin{align}
        F_{\mu_{n}}(t+\eps) + \eps \geq F_{\mu}(t) \quad \text{ for all } t \in \R \implies d(\mu_{n},\mu) \leq 2\eps.
    \end{align}
\end{proof}

\section{Expectation}

Expectation, in our new setting, corresponds to integration in a probability space. To this end, we start with simple functions.

\begin{definition}
    Let $(\Omega, \cF, \Prob)$ be a probability space. A measurable function $S: \Omega \to \R$ is called a \eax{simple random variable} if it takes finitely many values; that is, $S(\Omega) = \{s_{1},\ldots,s_{n}\}$ for some $n \in \N$.
\end{definition}
In this case, we can write $S = \sum_{i=1}^{n} c_{i} \1_{A_{i}}$ for some $c_{i} \in \R$ and $A_{i} \in \cF$. This representation need not be unique; however, note that any such simple function can be written in a unique way as
\begin{align}
    S = \sum_{i=1}^{n} \beta_{i} \1_{B_{i}}
\end{align}
where $\beta_{i} \in S(\Omega)$ are unequal and $B_{i}$'s are measurable sets that are mutually disjoint.

\begin{definition}
    The \eax{expectation of a simple random variable} $S = \sum_{i=1}^{n} c_{i} \1_{A_{i}}$ is defined as
    \begin{align}
        \E[S] \defeq \sum_{i=1}^{n} c_{i} \Prob(A_{i}).
    \end{align}
\end{definition}
Since this representation is not unique, one can verify that the expectation is well-defined by using the unique representation of $S$ as $\sum_{i=1}^{n} \beta_{i} \1_{B_{i}}$. Also, if $S$ and $T$ are simple random variables, then $S+T$ is also a simple random variable, and $\E[S+T] = \E[S] + \E[T]$. Similarly, for any $\alpha \in \R$, we have $\E[\alpha S] = \alpha \E[S]$. If $S \geq 0$, then $\E[S] \geq 0$. We can now extend the definition of expectation to all non-negative random variables, and then to all random variables.

\begin{definition}
    The \eax{expectation of a non-negative random variable} $X: \Omega \to [0,\infty)$ is defined as
    \begin{align}
        \E[X] \defeq \sup\{\E[S] \mid 0 \leq S \leq X, \, S \text{ is a simple random variable}\}.
    \end{align}
    If $X$ is a general random variable, we can write $X = X^{+} - X^{-}$, where $X^{+} = \max\{X,0\}$ and $X^{-} = \max\{-X,0\}$. If $\E[X^{+}] < \infty$ or $\E[X^{-}] < \infty$, then we define the \eax{expectation of (the) random variable} $X$ as
    \begin{align}
        \E[X] \defeq \E[X^{+}] - \E[X^{-}].
    \end{align}
\end{definition}

\begin{proposition}
    Let $X$ and $Y$ be non-negative random variables. Then, the following hold true.
    \begin{enumerate}
        \item $\E[X] \geq 0$.
        \item $\E[X+Y] = \E[X] + \E[Y]$, if at least one of $\E[X]$ or $\E[Y]$ is finite.
    \end{enumerate}
\end{proposition}
\begin{proof}
    The first part is clear from the definition since $S \equiv 0$ is a simple random variable such that $0 \leq S \leq X$, so $\E[X] \geq \E[S] = 0$. For the second part, let $S$ and $T$ be simple random variables such that $0 \leq S \leq X$ and $0 \leq T \leq Y$. Then, $S+T$ is a simple random variable such that $0 \leq S+T \leq X+Y$, and hence $\E[X+Y] \geq \E[S+T] = \E[S] + \E[T]$. Taking supremum over all such $S$ and $T$, we get $\E[X+Y] \geq \E[X] + \E[Y]$. For the reverse inequality, let $U$ be a simple random variable such that $0 \leq U \leq X+Y$. Then, we can write $U = U_{1} + U_{2}$ where $U_{1} = U\1_{\{X > 0\}}$ and $U_{2} = U\1_{\{Y > 0\}}$. Then, $0 \leq U_{1} \leq X$ and $0 \leq U_{2} \leq Y$, and hence $\E[U] = \E[U_{1}] + \E[U_{2}] \leq \E[X] + \E[Y]$. Taking supremum over all such $U$, we get $\E[X+Y] \leq \E[X] + \E[Y]$. Thus, we have equality.
\end{proof}

\subsection{Results of Convergence}

\noindent \textit{August 28th.}

We also have the monotone converge theorem, of which we state a basic version first.

\begin{theorem}
    Suppose $X \geq 0$ is a random variable, and $\{\varphi_{n}\}_{n=1}^{\infty}$ is a sequence of simple random variables such that $0 \leq \varphi_{n} \uparrow X$. Then, $\E[\varphi_{n}] \uparrow \E[X]$.
\end{theorem}
\begin{proof}
    Fix $\eps > 0$. By definition, there exists a simple function $S$ such that $0 \leq S \leq X$ and $\E[S] \geq \E[X] - \eps$. For a fixed $\omega \in \Omega$, there exists a natural $N_{\omega} \in \N$ such that $\varphi_{n}(\omega) \geq S(\omega)-\eps$ for all $n \geq N_{\omega}$. Now define $A_{n} = \{\omega \in \Omega \mid \varphi_{n}(\omega) \geq S(\omega)-\eps\}$; since $\varphi_{n}$'s are increasing functions, we have $A_{1} \subseteq A_{2} \subseteq \cdots \subseteq A_{n} \subseteq A_{n+1} \subseteq \cdots$, and $\bigcup_{n=1}^{\infty} A_{n} = \Omega$. Thus, given $\delta > 0$, we can find a natural $N_{\delta} \equiv N$ such that $\Prob(A_{N_{\delta}}) \geq 1-\delta$. Now, for all $n \geq N$, we have
    \begin{align}
        \E[\varphi_{n}] = \E[\varphi_{n} \1_{A_{N}}] + \E[\varphi_{n} \1_{A_{N}^{c}}].
    \end{align}
    But note that on such an $A_{N}$, we have $\varphi_{n} \geq S-\eps$, and hence $\E[\varphi_{n} \1_{A_{N}}] \geq \E[(S-\eps)\1_{A_{N}}] = \E[S\1_{A_{N}}] - \eps\Prob(A_{N})$. If $M = \sup_{\omega \in \Omega} S(\omega)$, then $\E[\varphi_{n} \1_{A_{N}^{c}}] \geq 0$ and $\E[S\1_{A_{N}}] = \E[S] - \E[S\1_{A_{N}^{c}}] \geq \E[S] - M\Prob(A_{N}^{c})$. Thus, we have
    \begin{align}
        \E[\varphi_{n}] \geq \E[S] - M\Prob(A_{N}^{c}) - \eps\Prob(A_{N}) \geq \E[S] - M\delta - \eps.
    \end{align}
    Choosing $\delta = \eps/M$, we have $\E[\varphi_{n}] \geq \E[S] - 2\eps \geq \E[X] - 3\eps$. Since $\eps > 0$ was arbitrary, we have $\lim_{n \toinf} \E[\varphi_{n}] = \E[X]$.
\end{proof}

\begin{theorem}[The \eax{monotone convergence theorem}]
    Suppose $X \geq 0$ is a random variable, and $\{X_{n}\}_{n=1}^{\infty}$ is a sequence of non-negative random variables such that $0 \leq X_{n} \uparrow X$. Then, $\E[X_{n}] \uparrow \E[X]$.
\end{theorem}

\begin{lemma}[\eax{Fatou's lemma}]
    Suppose $\{X_{n}\}_{n=1}^{\infty}$ is a sequence of non-negative random variables. Then,
    \begin{align}
        \E[\liminf_{n \to \infty} X_{n}] \leq \liminf_{n \to \infty} \E[X_{n}].
    \end{align}
\end{lemma}

\begin{proof}
    Define $Y_{n} = \inf_{k \geq n} X_{k}$ for all $n \in \N$. Then, $\{Y_{n}\}_{n=1}^{\infty}$ is a sequence of non-negative random variables such that $Y_{1} \leq Y_{2} \leq \cdots$ and $Y_{n} \uparrow \liminf_{n \toinf} X_{n}$. By the monotone convergence theorem, we have
    \begin{align}
        \E[\liminf_{n \toinf} X_{n}] = \lim_{n \toinf} \E[Y_{n}] \leq \liminf_{n \toinf} \E[X_{n}],
    \end{align}
    where the last inequality follows from the fact that $Y_{n} \leq X_{n}$ for all $n$.
\end{proof}

\textit{September 1st.}

\begin{theorem}[The \eax{dominated convergence theorem}]
    Suppose $\{X_{n}\}_{n=1}^{\infty}$ is a sequence of random variables such that $X_{n} \to X$ almost surely, and there exists a non-negative random variable $Y$ such that $\abs{X_{n}} \leq Y$ for all $n$, and $\E[Y] < \infty$. Then, $\E[X_{n}] \to \E[X]$.
\end{theorem}
\begin{proof}
    Define $Z_{n} = 2Y - \abs{X_{n}-X}$. Since $\Z_{n} \geq 0$, we can apply Fatou's lemma to get
    \begin{align}
        \E[\liminf_{n \toinf} Z_{n}] \leq \liminf_{n \toinf} \E[Z_{n}] \implies 2\E[Y] - \E[\limsup_{n \toinf} \abs{X_{n}-X}] \leq 2\E[Y] - \limsup_{n \toinf} \E[\abs{X_{n}-X}],
    \end{align}
    which implies that $\E[\limsup_{n \toinf} \abs{X_{n}-X}] \geq \limsup_{n \toinf} \E[\abs{X_{n}-X}]$. Now, since $X_{n} \to X$ almost surely, we have $\limsup_{n \toinf} \abs{X_{n}-X} = 0$ almost surely, and hence $\E[\limsup_{n \toinf} \abs{X_{n}-X}] = 0$. Thus, we have $\limsup_{n \toinf} \E[\abs{X_{n}-X}] = 0$, which implies that $\E[\abs{X_{n}-X}] \to 0$, and hence $\E[X_{n}] \to \E[X]$.
\end{proof}